\chapter[Discussion, Conclusion, and Future Work]{Discussion, Conclusion,\\and Future Work}
\label{chap:discussion}

\section{Discussion of Principal Findings}
Synthesizing the empirical results presented in Chapter~\ref{chap:results}, this thesis distills three key architectural lessons for multimodal deep survival analysis on modest-sized clinical cohorts:

First, \emph{the parameter-compression efficiency of CP decomposition must be evaluated within its specific architectural scope}. Replacing full outer-product projection layers with low-rank CP factorization reduces the parameter count of the three initial bilinear layers by \textbf{98.38\%} (from 229696 down to just 3712 parameters) and achieves a 90.90\% reduction across the complete interaction blocks including their downstream MLPs. However, this translates to an overall whole-network reduction of approximately 0.67\% because the 3D ResNet-18 volumetric encoder dominates the total parameter budget. This distinction clarifies that the primary benefit of applying CP decomposition in this architecture lies not in shrinking overall model storage size, but rather in \emph{preventing localized overfitting at the multimodal fusion bottleneck}---the exact location most vulnerable to parameter explosion when combining multiple feature streams on a training cohort of only 251 patients. The minor 0.0066 C-index trade-off between Full CP and Full Outer-Product represents a highly reasonable price for structural regularization at the interaction layer.

Second, \emph{activating more pairwise interactions does not automatically yield superior prognostic performance}. Across all seven surveyed CP architectures containing at least one pairwise interaction, the proposed RC-Free topology retains only the two MRI-centered interaction pairs (MR and MC) with four gated residual branches, yet achieves the highest mean validation C-index (\textbf{0.8347}). This result not only surpasses the three-pair Full CP model (0.8278) while saving 15460 parameters, but even the single-pair RC-only configuration reaches 0.8319, outperforming Full CP. The underlying reason is that all three input modalities already preserve their direct representation pathways through gated residual connections; consequently, forcing all three pairwise interactions simultaneously can introduce representational redundancy and increase estimation variance on a limited cohort. By positioning 3D structural MRI as the central hub connecting whole-brain radiomics and clinical scores, RC-Free provides the most balanced, compact, and effective fusion topology.

Third, \emph{the practical interpretation of risk stratification is inherently linked to sample size and cumulative event counts}. Kaplan--Meier survival analysis confirms that the log-risk scores generated by RC-Free separate patients into distinct progression trajectories, achieving log-rank $p=0.0134$ on the validation set (12/26 events in the high-risk group versus 4/28 in the low-risk group) and $p<0.001$ on the independent post-hoc test set. Nevertheless, when examining horizon-specific discrimination, the high validation AUCs at 2 and 3 years (0.9805 and 0.9788) rely on only \textbf{7 cumulative events}. By the 5-year horizon, the validation AUC decreases to 0.7889 with wider cross-seed standard deviation as the number of event-free subjects remaining under active follow-up beyond 5 years drops to 19 (Table~\ref{tab:horizon-support}).

\section{Study Limitations}
\label{sec:study-limitations}
Despite encouraging empirical results, this thesis acknowledges four methodological limitations:
\begin{enumerate}[itemsep=0.15em,topsep=0.25em]
\item \textbf{Data are currently restricted to the international ADNI cohort and have not yet been validated in Vietnamese hospitals.} All training, hyperparameter tuning, and testing in this study rely exclusively on the international ADNI database with modest evaluation split sizes (54 validation and 54 post-hoc test participants). The model has not yet been collected or externally validated on real-world patient cohorts from major hospitals in Vietnam, where MRI scanner protocols, clinical workflow practices, and demographic profiles may differ substantially.
\item \textbf{Retrospective time-axis design.} The survival time $t_i$ is measured from the third MRI scan date $p_i$, whereas the cohort inclusion criteria require participants to remain event-free through the 18-month milestone $L_i$ (creating potential immortal time bias if the model were interpreted as a prospective forecasting tool deployed immediately at scan date $p_i$ \cite{suissa2008immortal}). In addition, the $\pm90$-day clinical matching window results in 57 participants having an MMSE or CDR score linked to their third MRI scan that was actually recorded after date $p_i$ (Section~\ref{sec:data-split}), making the current findings strictly suited for retrospective evaluation on the stable post-$L_i$ cohort.
\item \textbf{Input predictors and outcome event criteria both rely on MMSE and CDR.} Although the input observation window and post-milestone outcome window are chronologically separated around the 18-month mark, defining progression events via MMSE and CDR thresholds inevitably couples the outcome definition to these two rating scales; the model has not yet been evaluated against independent specialist clinical diagnoses or pathological fluid/PET biomarkers.
\item \textbf{Model output is currently limited to relative risk ranking.} The scalar log-risk score $\eta_i$ from the Cox head establishes who is at higher risk of earlier progression relative to other patients, but has not yet been calibrated into an individualized absolute progression probability (in percentage terms at specific clinical horizons), nor has real-time inference latency been benchmarked on hospital workstation hardware.
\end{enumerate}

\section{Future Work}
Building directly upon the limitations identified above, I propose three practical directions for future research:
\begin{enumerate}[itemsep=0.15em,topsep=0.25em]
\item \emph{External validation on real-world hospital cohorts in Vietnam}: Collaborating with major domestic hospitals to collect longitudinal follow-up data from Vietnamese patients, evaluate external generalizability, adapt the model to local clinical settings, and incorporate readily available domestic laboratory markers.
\item \emph{Refining the temporal landmark prediction design}: Enforcing a strict cutoff ensuring that all input clinical assessments occur prior to the final MRI scan date $p_i$, or adopting a strict fixed-landmark survival design starting from milestone $L_i$ \cite{vanhouwelingen2007dynamic} to completely eliminate temporal overlap between input observations and outcome follow-up.
\item \emph{Calibrating individualized progression probabilities and enhancing clinical interpretability}: Estimating the baseline survival function via Breslow's estimator \cite{kvamme2019time} to convert relative log-risk scores into calibrated annual progression probabilities for individual patients, while analyzing learned attention weights and dynamic gating coefficients to provide transparent explanations and benchmarking inference runtime for clinical deployment.
\end{enumerate}

\section{Conclusion}
Prognosticating disease progression among individuals with mild cognitive impairment (MCI) requires not only tracking longitudinal disease trajectories across repeated clinical visits, but also harmonizing brain structural morphology, quantitative mathematical texture descriptors, and clinical cognitive assessments. Using a retrospective cohort of 359 participants from the ADNI initiative across three longitudinal visits, this thesis has developed an end-to-end multimodal deep survival analysis pipeline integrating 3D structural MRI volumes, whole-brain radiomics features, and clinical MMSE/CDR scores. Systematic ablation experiments confirm the contribution of each architectural design step: longitudinal MRI sequences modeled with T-LSTM and additive attention substantially outperform single-timepoint scans; integrating all three modalities markedly elevates risk-ranking concordance; and explicitly modeling pairwise bilinear interactions captures cross-modal relationships more effectively than standard feature concatenation.

Notably, the proposed \textbf{RC-Free} architecture combines two complementary simplification mechanisms: low-rank CP tensor decomposition reduces the parameter count of the three initial bilinear layers by 98.38\%, and pruning the direct RC interaction saves an additional 15460 parameters relative to Full CP while retaining all three input sources. On the validation set, RC-Free achieves the highest mean C-index among all seven pairwise CP topologies ($0.8347\pm0.0237$) and delivers statistically significant survival group separation on Kaplan--Meier curves ($p=0.0134$). Together, these findings demonstrate that combining low-rank tensor parameterization with domain-informed interaction pruning offers a compact, accurate, and promising paradigm for multimodal survival prognosis on sample-limited clinical datasets.
